\documentclass[12pt,twoside,a4paper]{report}

%====================== PACKAGES ======================
\usepackage{colortbl}
\usepackage[french]{babel}
\usepackage[utf8x]{inputenc}
%pour gérer les positionnement d'images
\usepackage{float}
\usepackage{amsmath}
\usepackage{graphicx}
\usepackage[colorinlistoftodos]{todonotes}
\usepackage{url}
%pour les informations sur un document compilé en PDF et les liens externes / internes
\usepackage{hyperref}
%pour la mise en page des tableaux
\usepackage{array}
\usepackage{tabularx}
%pour utiliser \floatbarrier
%\usepackage{placeins}
%\usepackage{floatrow}
%espacement entre les lignes
\usepackage{setspace}
%modifier la mise en page de l'abstract
\usepackage{abstract}
%police et mise en page (marges) du document
\usepackage[T1]{fontenc}
\usepackage[top=2cm, bottom=2cm, left=2cm, right=2cm]{geometry}
%Pour les galerie d'images
\usepackage{subfig}

%====================== INFORMATION ET REGLES ======================

%rajouter les numérotation pour les \paragraphe et \subparagraphe
\setcounter{secnumdepth}{4}
\setcounter{tocdepth}{4}

\hypersetup{                            % Information sur le document
pdfauthor = {Sow Mamadou,},          % Auteurs
pdftitle = {Rapport de stage -
            Développement web},           % Titre du document
pdfsubject = {Mémoire de Projet},       % Sujet
pdfkeywords = {Tag1, Tag2, Tag3, ...},  % Mots-clefs
pdfstartview={FitH}}                    % ajuste la page à la largueur de l'écran
%pdfcreator = {MikTeX},% Logiciel qui a crée le document
%pdfproducer = {}} % Société avec produit le logiciel

%======================== DEBUT DU DOCUMENT ========================

\begin{document}

%régler l'espacement entre les lignes
\newcommand{\HRule}{\rule{\linewidth}{0.5mm}}
\renewcommand{\thesection}{\arabic{section}}
%page de garde
\begin{titlepage}
\begin{center}

% Upper part of the page. The '~' is needed because only works if a paragraph has started.
\includegraphics[width=0.35\textwidth]{./logo_lyon1_2.jpg}~\\[0.2cm]

\textsc{\LARGE Université Claude Bernard Lyon 1}\\
Adresse: \textsc{43, Boulevard du 11 Novembre  \\ Villeurbanne 69100}\\[1cm]

\textsc{\Large }\\[0.5cm]

% Title
{\huge \bfseries Rapport De Stage\\}

\HRule \\[0.4cm]

{\huge \bfseries \Large Développement suivant un modèle MVC de La définition de permissions sur la console d'administration en fonction du profil des utilisateurs.} \\[0.4cm]
 %Definition des droits sur une console d'administration en fonction du profil des utilisateurs%

\HRule \\[1.5cm]

%logo entreprise
\includegraphics[width=0.35\textwidth]{./secuserve_logo.png}~\\[0.2cm]

%nom de l'entreprise
\textsc{\LARGE Entreprise Secuserve}\\
Adresse: \textsc{74, Bd du 11 Novembre 1918 \\ 69100 VILLEURBANNE}\\[2.5cm]

% Author and supervisor
\begin{minipage}{0.4\textwidth}
\begin{flushleft} \large
\emph{Stagiaire}\\
Nom: \textbf{\textsc{Sow}}\\
 Prénom: \textbf{\textsc{Mamadou}}\\
Numéro-étudiant: \textbf{\textsc{11414531}}\\
\textbf{\textsc{Licence 3 Informatique}}\\
Année-Scolaire: \textbf{\textsc{2016-2017}}
\end{flushleft}
\end{minipage}
\begin{minipage}{0.4\textwidth}
\begin{flushright} \large
 \textbf{\textsc{}}\\
\emph{Maitre de stage} \\
Nom: \textbf{\textsc{Fèvre  }}\\
Prénom: \textbf{\textsc{Pièrre   }}\\
\textbf{\textsc{Ingénieur DEVOPS}}\\
Lieu:   \textbf{\textsc{Lyon    }}\\
\end{flushright}
\end{minipage}




\vfill

% Bottom of the page
{\large \today}

\end{center}
\end{titlepage}

\tableofcontents
\thispagestyle{empty}
\setcounter{page}{0}
%ne pas numéroter le sommaire

%ne pas numéroter cette page
\thispagestyle{empty}
\newpage

\renewcommand{\abstractname}{Remerciements}
\addcontentsline{toc}{chapter}{Remerciements}
\begin{abstract}
\hskip7mm
\begin{spacing}{1.3}
Tout d'abord je tiens à remercier l'ensemble du personnel de l'entreprise Secuserve pour votre accueil chaleureux. \\
Au présidant Mr Stéphane Bouché de m'avoir donner la chance d'exploiter et d'acquérir des nouvelles expériences au seins de son entreprise.\\
Je remercie également mon maitre de stage Mr Pièrre Fèvre de m'avoir choisi et d'avoir eu confiance en moi pour intégrer dans cette structure jeune, dynamique, motivé et passionné de l'informatique.\\
Mais aussi à Mr Lionel Fèvre le responsable R\&D que je respecte énormément, qui regorge plein de ressource et de motivation avec lequel je travaille et apprend plein de nouvelles choses dans ce projet.\\

Enfin je remercie mon tuteur Fabien Duchateau pour le temps qu'il accorde à cette expérience professionnel ma référente pédagogique Marie Lefevre qui à toujours été présente quand j'en ai eu besoin durant toute cette année scolaire 2016-2017 à l'université Lyon 1.   
\end{spacing}
\end{abstract}
\thispagestyle{empty}



\newpage

%espacement entre les lignes d'un tableau
\renewcommand{\arraystretch}{1.5}

%====================== INCLUSION DES PARTIES ======================





\begin{center}
%\vspace{5cm }
\section{ \colorbox{orange}{Introduction}}
\end{center}
\begin{spacing}{1.3}
Stage en Licence 3 Informatique d’une durée de 4 mois en développement web suivant le modèle MVC (PHP, JAVA-SCRIPT, BOOTSTRAP, JQUERY, XML etc...), dans l'entreprise Secuserve qui est éditeur des services de sécurité et de messagerie.\\
Mon stage à débuté le 02/05/2017 et se terminera le 31/08/2017 dans les locaux des Secuserve à Villeurbanne, Au-delà d'enrichir et de mettre en pratique mes connaissances acquises, j'ai pu confirmer mon souhait de devenir développeur web.\\ \\
L'objet du stage porte sur deux aspects en corrélation:
D'une part la définition de permissions sur la console d'administration en fonction du profil des utilisateurs issu d'une structure hiérarchique.\\
D'autre part la définition des règles de communication et de sécurité pour le filtrage des messages en sortie.\\
Les développements suivront le modèle mvc et pourront utiliser des outils tel-que: java script, jquery, bootstrap, etc..., tout développement devra être testé et documenté.\\
Travaillant avec le responsable de R\&D , j'apprends dans d'excellentes conditions en adoptant des méthodes de travail pratiques et efficaces.  \\

En vue de rendre compte de manière fidèle et analytique le temps passé au sein de cette entreprise depuis le mois de mai jusqu'à présent, je vais commencer par vous présenter l'environnement professionnel du stage en suite tous les aspects de mon travail jusqu'à maintenant et faire une mini présentation de la seconde partie de mon stage qui commence après le mois de juin en précisant les différentes missions et tâches que j'ai pu effectuer au sein du service. Je finirai par la conclusion  en résumant les apports de mon stage.
\end{spacing}
\newpage

\section{Environnement Professionnel }
\subsection{Presentation de secuserve}
\subsubsection{Histoire}
Fondée en 2003, Secuserve est la première société française à éditer des services de sécurité et de messagerie, disponibles à la demande via Internet.\\
Secuserve est reconnue pour son expertise dans le domaine du SaaS, de la sécurité, de la collaboration et de la messagerie, son excellence dans l’accompagnement des projets d’externalisation de ses clients et la qualité de ses services.\\
Secuserve est basé dans deux endroits suivant: 
\begin{itemize}
\item un siège social à Paris 
\item une Agence à Lyon
\end{itemize}

\subsubsection{Organisation de l'entreprise}
Comme évoqué précédemment, cette entreprise est une petite structure dont le siège social est à Paris mais qui possède des bureaux à Lyon (Villeurbanne pour être précis). 
Secuserve évolue sur le marché global de l’informatique mais est spécialisée sur le segment de la sécurité lié à la messagerie électronique.
\paragraph{Organigramme Secuserve}
Le chef d’entreprise Monsieur Stéphane Bouche l’a créée en 2003 suite à la montée des cybers attaques via les comptes de messagerie dans les entreprises et administrations. Fort de son expérience, Secuserve propose donc des services de protection des emails entrant et sortant à toutes les entreprises quel que soit leur domaine d’activité ou leur taille. Si cela résume ce qu’est l’offre Secuserve, celle-ci se décline en trois produits ou solutions distincte.
Concrètement, la structure est répartie sur les deux zones (Paris et Villeurbanne).\\
Les équipes de Paris sont dédiées au commerce avec trois commerciaux dont la directrice commerciale Mme Aline Breton. Le directeur technique est également installé dans la capitale où sont les datacenters de la société pour une question de proximité en cas de problème. Le site de Lyon est dédié à la recherche et développement dont l’activité est reconnue par le ministère de l’enseignement et de la recherche. J’ai été recruté pour développer la partie commerciale sur le sud de la France avec Monsieur Edouard Courbon, ingénieur d’affaires,  qui a quitté la société en Février 2017 pour un post aux Etats-Unis. De fait, nous serons deux sur cette partie à compter de fin février 2016 : le PDG et moi.
Tout ce qui appartient au domaine administratif et comptable est géré par Mme Gaillac à Paris.

\subsubsection{L’offre Secuserve}
Secuserve propose trois services, qui sont les suivants:
\paragraph{E-securemail}
Il s’agit d’un service de filtrage (antispam et antivirus) et de sécurisation des emails.
\paragraph{Optimails:} C’est un service de messagerie d’entreprise collaborative et externalisée.
\paragraph{Sharecan:}Service de stockage et de partage de fichiers en ligne.
Il s’agit d’un marché très spécifique mais sur lequel la concurrence est accrue.

\paragraph{concurrence:}
Il s’agit d’un marché très spécifique mais sur lequel la concurrence est accrue.\\

%tableau centré à taille variable qui s'ajuste automatiquement suivant la longueur du contenu
\begin{figure}[!h]
\begin{center}
\begin{tabular}{|l|l|l|l|c|l|l|}
  \hline
 Entreprise & Chiffre d 'affaire 2015& Parts marché\\
  \hline
   Vade Secure& 4 227 220 euros & 28.49\% \\
   \hline
  Mail in Black& 2 695 500 euros & 18.17\% \\
  \hline
  AltoSpam& 1 038 700euros & 7.00\% \\
  \hline  
   \rowcolor{orange}Secuserve& 873 700 euros& 5.85\% \\
	\hline
   Autres(microsoft,symentec,etc...)& 6 000 000 euros & 40\%\\
	\hline  
    Total& 14 835 100 & 100\% \\
  \hline
\end{tabular}
\end{center}
\caption{Parts de marché du secteur}
\end{figure}


\newpage

%mission
\section{Missions}

\subsection{Missions réalisées}
\subsubsection{Mise en place de l'environnent de travail}
Pour pouvoir faire ces taches de développement il fallait avoir certains notions de la structure des fichiers et quelque logiciels de bases a installer pour pouvoir débuter le développement du site appeler (la console d'administration) tel-que: 
\begin{enumerate}
\item la structure de la base de donnée \textbf{Mysql} qui comprend plusieurs serveurs de bases de données ces derniers correspondent à des environnements de travail pour les développeurs, qui sont différentes étapes du développement (développement, pré-production, production) avant la version final.
\item la structure du dépôt de fichiers gérer par le gestionnaire de versions \textbf{SVN} qui nous permets de faire des commits sur le serveur de développement.
\item Installer les IDE \textbf{(NETBEANS et PhpStorm)} car le développement actuelle est fait sur \textbf{Netbeans} mais par la suite on développera avec \textbf{PhpStorm} et qu'on utilisera \textbf{Gitlab} comme gestionnaire de version. Il fallait aussi installer \textbf{PUTTY} pour la connexion aux serveur du développement et \textbf{XMIND} un logiciel qui permet de creer des cartes en \textbf{XML}. 
\end{enumerate}
  
\subsubsection{Définition de permissions sur la console d'administration en fonction du profil des utilisateurs.}
\begin{description}


\item[Profile:] un profil est une configuration de droits (admin, rssi etc..)

\item[Login:] utilisateur de la console d’administration ayant un profil.

\item[Pages mvcisé:] on à utiliser le modèle view contrôleur pour développer les différents aspects de notre site.

\item[Pages non mvcisé:] utilises un contrôleur pour contrôler les droits afin de rediriger vers la page correspondante en attendant que la page soit mvcisé.

\item[Contrôleur:] Permet de gérer définir des actions à travers ses méthodes sur un aspect du site web.

\item[Méthode:] Permet de gérer les actions à réaliser avant d'envoyer les données sur la vue (page web).
\end{description}

\paragraph{Gestion des droits}
Est une interface d'administration des droits que j'ai développé,
Comme toutes les autres pages mvcisé j'utilise un contrôleur qui hérite d'un contrôleur générique (créer par responsable R\&D) dans lequel on défini une méthode setRights qui permet:
\begin{enumerate}


\item Récupérer dans une variable les noms de tous les contrôleurs et leurs méthodes en parcourant un fichier XML générer à partir d'une carte XMIND.
\item Ensuite créer un formulaire (tableau avec des checkbox) d'administration à partir de cette variable pour les différents profils.
	\begin{itemize}
	\item Lorsqu'on fini de cocher les cases du formulaire et qu'on valide pour l'envoyer, l'action est automatiquement historisée en enregistrant les données modifiées suivi des messages de session en cas d'erreur ou de succès.\label{annexe1}annex
	\end{itemize}
\end{enumerate}
\paragraph{Autorisations}
\subparagraph{Controlleur: generiqueController} 
Dans chaque contrôleur on vérifie les droits d’utilisateur de tous les liens sur la pages avant de l’afficher, en parcourant
un fichier XML qui définit l’architecture de celle-ci.
\subparagraph{Exception pour superuser}
Par défaut si on est super-user le contrôleur donne droit sur toutes les pages du site.
\subparagraph{Pages publiques} 
Certaine page sont public comme par exemple les pages de l'aide (documentation, téléchargement, etc...) et son
accessible sans vérification de droits.

\subparagraph{Traduction des pages}
La console d'administration utilise deux formats de langue, par conséquent il faut chercher chaque mot dans la table correspondance (développer par le responsable D\&R Lionel Fevre) pour ensuite l'appeler dans une fonction qui permet d'afficher ce dernier dans la langue choisie(FR,EN).

\subsection{Missions à réalisées}

\subsection{Autres Missions}





%\input{./existant.tex}
 
%\input{./besoins.tex}

%\input{./autre_partie.tex}

%\input{./resultats.tex}

%\input{./bilan.tex}

%\input{./annexes.tex}

\newpage

%récupérer les citation avec "/footnotemark"
\nocite{*}

%choix du style de la biblio
\bibliographystyle{plain}
%inclusion de la biblio
%\bibliography{bibliographie.bib}
%voir wiki pour plus d'information sur la syntaxe des entrées d'une bibliographie

\end{document}